
This paper proposes a confidence geometry framework for detecting non-terminating proof searches in neural theorem proving systems. The central claim is that LLM confidence trajectories encode geometric structure of successful proof spaces, and measuring entropy variance over proof search trajectories can signal probable non-termination. The research presents results from 100 experiments purported to validate this principle, reporting strong correlation between confidence variance and timeout outcomes (r=0.80, p<10⁻²³). However, critical examination of the research artifacts reveals that the experimental results are based on synthetically generated mock data rather than actual LeanDojo theorem proving experiments. The results files contain entries labeled "mock_theorem_XXX" with variance values generated by predetermined statistical distributions. While the conceptual framework and methodology are described in detail, the empirical validation remains incomplete. The paper acknowledges that Phase 5 baseline comparison was not completed, but does not disclose that the reported Phase 4 results derive from synthetic rather than real data. This rewrite presents the research design and claims while clearly distinguishing between proposed methodology and actual empirical validation.
